\section{Transport under physical blocking}\label{sec:blocking_transport}

For a tensor $A=(A^i)_{i=0}^{d-1}$ and a block length $p$, write
$A^{[p]}$ for the tensor whose physical matrices are the consecutive products
$A^{i_1}\cdots A^{i_p}$. We keep the number $N$ of blocked sites distinct from
the number $Np$ of original sites. The exact Hilbert-space transport identities
below are elementary consequences of this definition and of the local-space
construction; they are not attributed to a source theorem.

\begin{definition}[Blocked configuration isometry]
    \label{def:blocked_configuration_isometry}
    \lean{MPSTensor.blockedConfigLinearIsometryEquiv}
    \lean{MPSTensor.blockedConfigLinearIsometryEquiv_apply_apply}
    \lean{MPSTensor.blockedConfigLinearIsometryEquiv_symm_apply_apply}
    \leanok
    Grouping each consecutive $p$-tuple of physical indices defines a canonical
    computational-basis unitary satisfying
    \begin{align}
        W_{N,p} & :\ell^2\!\left(\Fin{d^p}^{N}\right)
        \longrightarrow\ell^2\!\left(\Fin{d}^{Np}\right).
        \label{eq:blocked_configuration_isometry}
    \end{align}
    Its forward and inverse coordinate formulas are the corresponding
    reindexings by the blocked-configuration equivalence and its inverse.
\end{definition}

\begin{lemma}[Boundary-map transport under blocking]
    \label{lem:blocked_boundary_map_transport}
    \lean{MPSTensor.groundSpaceMap_blockTensor_blockedConfigEquiv}
    \lean{MPSTensor.blockedConfigLinearIsometryEquiv_groundSpaceMapES}
    \leanok
    \uses{def:blocked_configuration_isometry, def:blocked_tensor, def:ground_space_map, def:ground_space_map_es}
    For every $x\in\ell^2(\Fin{D}\times\Fin{D})$, the Euclidean boundary maps
    intertwine with the blocked-configuration isometry:
    \begin{align}
        W_{N,p}\Gamma_N^{\mathrm{ES},A^{[p]}}(x)
         & =\Gamma_{Np}^{\mathrm{ES},A}(x).
        \label{eq:blocked_boundary_map_transport}
    \end{align}
\end{lemma}

\begin{proof}
    \leanok
    \uses{lem:eval_word_block, lem:ground_space_map_apply, thm:ground_space_map_es_frobenius}
    For a blocked configuration $J=(J_1,\ldots,J_N)$, its flattened word
    $\sigma$, and a boundary matrix $X\in\MN{D}$, blocked-word evaluation gives
    \begin{align}
        (A^{[p]})^{J}
         & =A^{\sigma_1}\cdots A^{\sigma_{Np}},
        \label{eq:blocked_word_flattening}                        \\
        \Gamma_N^{A^{[p]}}(X)(J)
         & =\tr\!\left(A^{\sigma_1}\cdots A^{\sigma_{Np}}X\right)
        =\Gamma_{Np}^{A}(X)(\sigma).
        \label{eq:blocked_boundary_map_transport_coefficients}
    \end{align}
    Applying the Frobenius transport formula to the coefficient identity yields
    (\ref{eq:blocked_boundary_map_transport}).
\end{proof}

\begin{lemma}[Local spaces under coisometric reconstruction]
    \label{lem:coisometric_reconstruction_local_space}
    \lean{MPSTensor.groundSpace_eq_of_coisometry_reconstruction}
    \lean{MPSTensor.CPSVCanonicalFormData.%
        groundSpace_blockTensor_eq_iSup_representatives}
    \leanok
    \uses{def:ground_space_parent, lem:block_sum_local_ground_space}
    Let $A=(A^i)$ and $C=(C^i)$ have the same physical alphabet, and suppose
    there is a coisometry $U:\C^D\to\C^n$ such that
    $UU^\dagger=I_n$ and $A^i=U^\dagger C^iU$ for every $i$. Then, for every
    $L\geq1$,
    \begin{align}
        G_L(A) & =G_L(C).
        \label{eq:coisometric_reconstruction_local_space}
    \end{align}
    In particular, if $(A_j)_{j=1}^g$ are the distinct BNT representatives
    selected from literal canonical-form data for $A$, then, for $p,L\geq1$,
    \begin{align}
        G_L\!\left(A^{[p]}\right)
         & =\bigvee_{j=1}^{g}G_L\!\left(A_j^{[p]}\right).
        \label{eq:cpsv_blocked_representative_local_space}
    \end{align}
    The first equality is a derived trace identity for an exact coisometric
    reconstruction. It is not the literal direct-sum canonical-form equation
    of CPSV16.
\end{lemma}

\begin{proof}
    \leanok
    \uses{lem:block_sum_local_ground_space}
    For every word $w$ of positive length,
    $A^w=U^\dagger C^wU$. Cyclicity of the trace gives
    \begin{align}
        \Gamma_L^A(X)
         & =\Gamma_L^C(UXU^\dagger),
        \notag                        \\
        \Gamma_L^C(Y)
         & =\Gamma_L^A(U^\dagger YU),
        \notag
    \end{align}
    where the second equality uses $UU^\dagger=I_n$. This proves
    (\ref{eq:coisometric_reconstruction_local_space}). For literal
    canonical-form data, apply this equality to the retained direct sum.
    Lemma~\ref{lem:block_sum_local_ground_space} gives the sum over all
    displayed normal blocks. Blocks in one BNT phase class differ by a
    nonzero phase and an invertible gauge, so their local spaces agree; the
    sum therefore reduces to the distinct representatives, which proves
    (\ref{eq:cpsv_blocked_representative_local_space}).
\end{proof}

\begin{theorem}[A simultaneous one-letter span gives nearest-neighbor ground spaces]
    \label{thm:block_diagonal_one_site_span_parent_spanning}
    \lean{MPSTensor.WordTupleSpanTop.of_one_of_pos}
    \lean{MPSTensor.WordTupleSpanTop.exists_one_letter_identity_coefficients}
    \lean{MPSTensor.%
        pgvwc07_mem_iSup_groundSpace_of_trace_decomposition_of_identity_coefficients}
    \lean{MPSTensor.%
        pgvwc07_mem_iSup_groundSpace_of_iSup_restrictions_of_identity_coefficients}
    \lean{MPSTensor.%
        pgvwc07_iSup_groundSpace_eq_restriction_intersection_of_identity_coefficients}
    \lean{MPSTensor.%
        chainGroundSpace_toTensorFromBlocks_le_iSup_groundSpace_of_wordTupleSpanTop_one}
    \lean{MPSTensor.%
        exists_blockDiagonal_boundary_of_chainGroundSpace_of_wordTupleSpanTop_one}
    \lean{MPSTensor.%
        exists_blockDiagonal_boundary_chainGroundSpace_of_wordTupleSpanTop_one}
    \lean{MPSTensor.%
        chainGroundSpace_toTensorFromBlocks_eq_iSup_of_wordTupleSpanTop_one}
    \lean{MPSTensor.%
        hasParentHamiltonianGroundSpaceSpanning_toTensorFromBlocks_of_wordTupleSpanTop_one}
    \leanok
    \uses{lem:block_diagonal_chain_space_comparison_parent_spanning,
        thm:chain_ground_space_eq_mpv_blk}
    Let $B=\bigoplus_{j=1}^{g}\mu_jA_j$, where every $\mu_j$ is nonzero and
    every bond dimension is positive. Suppose that the simultaneous one-letter
    evaluations span the product algebra:
    \begin{align}
        \spn_{\C}\left\{(A_1^i,\ldots,A_g^i):0\leq i<d\right\}
         & =\prod_{j=1}^{g}\MN{D_j}.
        \label{eq:block_diagonal_one_site_span}
    \end{align}
    Then, for every $N>2$,
    \begin{align}
        \mathcal G_{N,2}(B)
         & =\bigvee_{j=1}^{g}\mathcal G_{N,2}(A_j),
        \notag                                                   \\
        \ker H_2^{(N)}(B)
         & =\spn\left\{\ket{V^{(N)}(A_j)}:1\leq j\leq g\right\}.
        \label{eq:block_diagonal_one_site_span_parent_spanning}
    \end{align}
\end{theorem}

\begin{proof}
    \leanok
    \uses{lem:block_diagonal_chain_space_comparison_parent_spanning,
        thm:chain_ground_space_eq_mpv_blk}
    Equation~(\ref{eq:block_diagonal_one_site_span}) supplies common
    coefficients $c_i$ satisfying $\sum_i c_iA_j^i=I_{D_j}$ for every $j$.
    In the restriction-intersection argument of the proof of
    \cite[Theorem~12]{PerezGarcia2007Matrix}, these coefficients replace the
    right-canonical identity and yield the first equality. This conditional
    variant and its relation to the source proof are documented in
    \cite{gap:pgvwc07_common_identity_coefficients}. The same span makes every
    $A_j$ injective, so
    Theorem~\ref{thm:chain_ground_space_eq_mpv_blk} identifies its periodic
    chain space with $\spn\{\ket{V^{(N)}(A_j)}\}$. The block-diagonal
    comparison of
    Lemma~\ref{lem:block_diagonal_chain_space_comparison_parent_spanning}
    then gives (\ref{eq:block_diagonal_one_site_span_parent_spanning}).
\end{proof}

\begin{lemma}[The spanning equation depends only on the local space]
    \label{lem:parent_spanning_of_local_space_eq}
    \lean{MPSTensor.chainGroundSpace_eq_of_groundSpace_eq}
    \lean{MPSTensor.HasParentHamiltonianGroundSpaceSpanning.%
        of_groundSpace_eq}
    \leanok
    \uses{def:parent_hamiltonian_ground_space_spanning,
        thm:parent_hamiltonian_kernel_eq_chain_ground_space}
    Let $B$ and $C$ have the same physical alphabet and satisfy $G_L(B)=G_L(C)$.
    Their periodic chain spaces agree for every $N$:
    \begin{align}
        \mathcal G_{N,L}(B)
         & =\mathcal G_{N,L}(C).
        \label{eq:parent_spanning_local_space_chain_congruence}
    \end{align}
    Their parent-Hamiltonian kernels agree for every $N>L$:
    \begin{align}
        \ker H_L^{(N)}(B)
         & =\ker H_L^{(N)}(C).
        \label{eq:parent_spanning_local_space_kernel_congruence}
    \end{align}
    Consequently, any fixed BNT-vector family satisfying the parent-Hamiltonian
    spanning equation for $C$ also satisfies it for $B$.
\end{lemma}

\begin{proof}
    \leanok
    \uses{thm:parent_hamiltonian_kernel_eq_chain_ground_space}
    When $N=0$ or $L>N$, both chain spaces are $\top$ by definition. Otherwise
    $0<N$ and $L\leq N$. Every translated local constraint is then the inverse
    image of the same subspace $G_L(B)=G_L(C)$, so their intersections are
    equal. This gives
    (\ref{eq:parent_spanning_local_space_chain_congruence}) for every $N$.
    When $N>L$,
    Theorem~\ref{thm:parent_hamiltonian_kernel_eq_chain_ground_space} identifies
    those intersections with the two parent-Hamiltonian kernels and gives
    (\ref{eq:parent_spanning_local_space_kernel_congruence}).
\end{proof}

\begin{theorem}[Original-lattice parent-Hamiltonian range]
    \label{thm:cpsv_original_lattice_parent_hamiltonian_range}
    \lean{MPSTensor.IsCPSVCanonicalForm.%
        exists_bnt_hasParentHamiltonianGroundSpaceSpanning}
    \leanok
    \uses{def:cpsv_canonical_form, def:bnt_cpsv,
        def:parent_hamiltonian_ground_space_spanning}
    Let $A$ be in literal CPSV canonical form with positive ambient bond
    dimension $D$, and suppose that the physical dimension satisfies $1<d$.
    There are a basis of normal tensors $(A_j)_{j=1}^{g}$ and a positive
    integer $L$ such that
    \begin{align}
        D^2<d^L,
        \qquad
        L\leq3D^5,
        \label{eq:cpsv_original_lattice_parent_hamiltonian_dimension_range}
    \end{align}
    and, for every original periodic chain of length $N>L$,
    \begin{align}
        \ker H_L^{(N)}(A)
         & =\spn\left\{\ket{V^{(N)}(A_j)}:1\leq j\leq g\right\}.
        \label{eq:cpsv_original_lattice_parent_hamiltonian_range}
    \end{align}
    Thus the orthogonal complement of the local MPS space is nonzero, and
    the corresponding bounded-range interaction is a parent Hamiltonian.
    This is the construction of
    \cite[Definition~3.9 and line~527]{Cirac2016MPDO_arXiv}, with its
    condition $d^L>D^2$ made explicit.
\end{theorem}

\begin{proof}
    \leanok
    \uses{thm:cpsv_canonical_form_bnt_refinement,
        thm:cpsv_bnt_dimension_bound,
        thm:gs_eq_bnt_span_source_normalization,
        lem:cyclic_window_range_antitonicity,
        lem:bnt_span_le_parent_hamiltonian_kernel,
        lem:coisometric_reconstruction_local_space,
        lem:parent_spanning_of_local_space_eq,
        thm:exact_pow_four_block_injectivity_normal_left_canonical,
        thm:chain_ground_space_eq_mpv_normal,
        thm:block_diagonal_one_site_span_parent_spanning,
        thm:unital_gauge_of_pd_fixed_point,
        thm:pgvwc07_dual_diag_unital}
    Choose one normal representative from each BNT phase class and put each
    representative in a Perron gauge for which
    \begin{align}
        \sum_a A^j_a A^{j\,\dagger}_a          & =\Id,
        &
        \sum_a A^{j\,\dagger}_a\Lambda_j A^j_a & =\Lambda_j
        \notag
    \end{align}
    with $\Lambda_j$ positive definite. Their dimensions satisfy
    $\sum_jD_j\leq D$, so $g\leq D$, and quantum Wielandt makes every
    representative block-injective at the common length $L_0=D^4$.

    Suppose first that $g\geq2$. Put
    \begin{align}
        R & =3(g-1)(L_0+1)+1,
        &
        L_{\mathrm{aux}} & =3(g-1)(L_0+1)+L_0.
        \notag
    \end{align}
    If $N>L_{\mathrm{aux}}$, then $N\geq R+L_0$. The source-normalized
    form of \cite[Theorem~12]{PerezGarcia2007Matrix} therefore identifies
    the range-$R$ parent-Hamiltonian kernel with the BNT span. Since
    $R\leq L_{\mathrm{aux}}$, cyclic-window antitonicity places the
    range-$L_{\mathrm{aux}}$ kernel inside the range-$R$ kernel. Every BNT
    vector satisfies the range-$L_{\mathrm{aux}}$ local constraints, which
    gives the reverse inclusion. Thus the range-$L_{\mathrm{aux}}$ kernel
    is exactly the BNT span for every $N>L_{\mathrm{aux}}$, and
    \begin{align}
        L_{\mathrm{aux}}
         & \leq 3(D-1)(D^4+1)+D^4
          \leq 3D^5.
        \notag
    \end{align}
    If $g=1$, the same spanning conclusion holds at
    $L_{\mathrm{aux}}=D^4+1$ by the normal-tensor intersection and closure
    property reviewed in \cite[Section~IV.C]{Cirac2021Matrix}. If $g=0$,
    the one-site block-diagonal spanning result gives it at
    $L_{\mathrm{aux}}=2$; this is a formal boundary extension of the source
    argument.

    In every case enlarge the interaction range to the common value
    $L=3D^5$. Range antitonicity gives the upper inclusion at $L$, while the
    usual frustration-free inclusion gives the reverse one, so the
    all-$N>L$ spanning equation is preserved. Since $D>0$ and $d\geq2$,
    \begin{align}
        D^2<2^{2D}\leq d^{2D}\leq d^{3D^5}=d^L.
        \notag
    \end{align}
    Exact coisometric reconstruction identifies the local space of the
    representative direct sum with $G_L(A)$ and transfers the equality to
    the original tensor. The chosen range is positive and satisfies
    $L\leq3D^5$ by equality.
\end{proof}

\begin{theorem}[Sharp blocking gives a nearest-neighbor parent Hamiltonian]
    \label{thm:cpsv_blocked_nearest_neighbor_parent_hamiltonian}
    \lean{MPSTensor.IsCPSVCanonicalForm.%
        exists_bnt_blocked_hasNearestNeighborParentHamiltonian}
    \leanok
    \uses{def:cpsv_canonical_form, def:bnt_cpsv, def:blocked_tensor,
        def:parent_hamiltonian_ground_space_spanning}
    Let $A$ be in literal CPSV canonical form with positive ambient bond
    dimension $D$. There are a basis of normal tensors
    $(A_j)_{j=1}^{g}$ and an integer $p$ satisfying
    $1\leq p\leq 3D^5$ such that $(A_j^{[p]})_{j=1}^{g}$ is a basis of
    normal tensors for $A^{[p]}$ and, for every $N>2$,
    \begin{align}
        \ker H_2^{(N)}\!\left(A^{[p]}\right)
         & =\spn\left\{\ket{V^{(N)}\!\left(A_j^{[p]}\right)}:
        1\leq j\leq g\right\}.
        \label{eq:cpsv_blocked_nearest_neighbor_parent_hamiltonian}
    \end{align}
    Thus the actual blocked tensor has a nearest-neighbor parent Hamiltonian
    on the blocked lattice. This is distinct from
    Theorem~\ref{thm:cpsv_original_lattice_parent_hamiltonian_range}, which
    supplies a nonzero bounded-range interaction and the parent-Hamiltonian
    spanning identity on every sufficiently long original chain. The two
    conclusions and the argument closing the short-ring window are recorded in
    \cite{gap:cpsv16_parent_hamiltonian_range_short_ring}.
\end{theorem}

\begin{proof}
    \leanok
    \uses{thm:cpsv_canonical_form_bnt_refinement,
        thm:cpsv_bnt_dimension_bound, thm:cpsv_bnt_sharp_span_ambient_cap,
        thm:cpsv_bnt_blocked_simultaneous_span,
        thm:cpsv_bnt_positive_blocking,
        lem:coisometric_reconstruction_local_space,
        lem:block_sum_local_ground_space,
        thm:block_diagonal_one_site_span_parent_spanning,
        lem:parent_spanning_of_local_space_eq}
    Choose $p$ so that the simultaneous one-letter evaluations of the blocked
    normal representatives span the product of their full matrix algebras.
    Positive blocking preserves the BNT property. Set
    $C^i=\bigoplus_{j=1}^{g}(A_j^{[p]})^i$. Then
    Theorem~\ref{thm:block_diagonal_one_site_span_parent_spanning} gives, for
    every $N>2$,
    \begin{align}
        \ker H_2^{(N)}(C)
         & =\spn\left\{\ket{V^{(N)}(A_j^{[p]})}:1\leq j\leq g\right\}.
        \label{eq:cpsv_blocked_representative_parent_spanning}
    \end{align}
    After the repeated phase-gauge copies and the ambient coisometry are
    removed, Lemma~\ref{lem:coisometric_reconstruction_local_space} gives
    \begin{align}
        G_2(A^{[p]})
         & =\bigvee_{j=1}^{g}G_2(A_j^{[p]})=G_2(C).
        \label{eq:cpsv_actual_direct_sum_local_space}
    \end{align}
    Lemma~\ref{lem:parent_spanning_of_local_space_eq} identifies the two
    parent-Hamiltonian kernels, so
    (\ref{eq:cpsv_blocked_representative_parent_spanning}) becomes
    (\ref{eq:cpsv_blocked_nearest_neighbor_parent_hamiltonian}).
\end{proof}

\begin{theorem}[Local ground-space transport under blocking]
    \label{thm:blocked_local_ground_space_transport}
    \lean{MPSTensor.groundSpaceES_blockTensor_map}
    \lean{MPSTensor.starProjection_groundSpaceES_blockTensor_map_apply}
    \lean{MPSTensor.starProjection_groundSpaceES_blockTensor_conj}
    \lean{MPSTensor.starProjection_groundSpaceES_three_blockTensor_conj}
    \leanok
    \uses{def:blocked_configuration_isometry, def:blocked_tensor, def:ground_space_es, lem:blocked_boundary_map_transport}
    The Euclidean local space of the blocked tensor is carried exactly to the
    original Euclidean local space:
    \begin{align}
        W_{N,p}G_N^{\mathrm{ES}}(A^{[p]})
         & =G_{Np}^{\mathrm{ES}}(A).
        \label{eq:blocked_local_ground_space_transport}
    \end{align}
    Hence their orthogonal projectors satisfy
    \begin{align}
        P_{G_{Np}^{\mathrm{ES}}(A)}
         & =W_{N,p}P_{G_N^{\mathrm{ES}}(A^{[p]})}W_{N,p}^{*}.
        \label{eq:blocked_ground_projector_conjugacy}
    \end{align}
    These statements include $N=0$ or $p=0$; in those cases the same explicit
    configuration isometry and boundary-map calculation apply. In particular,
    the projector identity holds for three blocked sites and $3p$ original
    sites. These identities concern the full local ground spaces and their
    orthogonal projectors under blocking.
    % Scope: this theorem does not identify the left or tail open-chain spectator
    % spaces needed by a three-block estimate and does not compare the blocked and
    % original parent Hamiltonians.
\end{theorem}

\begin{proof}
    \leanok
    \uses{lem:blocked_boundary_map_transport, thm:range_ground_space_map_es}
    The exact range identification for the Euclidean boundary map and
    (\ref{eq:blocked_boundary_map_transport}) give
    \begin{align}
        W_{N,p}\operatorname{range}\!\left(\Gamma_N^{\mathrm{ES},A^{[p]}}\right)
         & =\operatorname{range}\!\left(\Gamma_{Np}^{\mathrm{ES},A}\right),
        \label{eq:blocked_boundary_map_range_transport}
    \end{align}
    which is (\ref{eq:blocked_local_ground_space_transport}). For every
    subspace $K$, orthogonal projection under the unitary $W_{N,p}$ satisfies
    \begin{align}
        P_{W_{N,p}(K)} & =W_{N,p}P_KW_{N,p}^{*}.
        \label{eq:projection_under_blocked_isometry}
    \end{align}
    Taking $K=G_N^{\mathrm{ES}}(A^{[p]})$ gives
    (\ref{eq:blocked_ground_projector_conjugacy}).
\end{proof}
